\chapter{Introducción}
\label{ch:no_lineal}

En este capítulo se presenta el contexto del trabajo, así como la motivación del problema concreto que se aborda. Se explica los objetivos y sus limitaciones, además de las herramientas de trabajo utilizadas. Por último, se explica el esquema general de la memoria del proyecto.

\section{Contexto de trabajo}

% TODO: Explicar en qué trabaja la empresa, cómo trabaja, etc.

 Ibernex está especializada en el diseño, desarrollo e integración de soluciones y servicios tecnológicos destinados al sector socio-sanitario.
 Realizan soluciones para automatizar y digitalizar la atención y experiencia de residencias u hospitales. \\


 La compañía realiza todo el proceso de construcción del producto ya que se encarga de la fabricación del hardware y de la implementación del software de las soluciones para cada cliente.
 Los clientes de la empresa, como se ha comentado anteriormente, son residencias u hospitales que quieren digitalizar el proceso de cuidado y atención de pacientes, además de otros procesos que puedan tener según sus necesidades. \newline

% En este caso interesa la parte del software. Ibernex trabaja con una aplicación 

Duda - TODO: explicar qué es lo que tiene implementado Ibernex para sus soluciones y qué me presenta como punto de partida para desarrollar el sistema ???

% \begin{figure}[!h]
%     \centering
%     \includegraphics[width=0.8\textwidth,height=6cm]{Imagenes/Arte_abstracto}
%     \caption{Un pie de foto}
%     \label{fig:una_etiqueta}
% \end{figure}




\section{Objetivos y limitaciones}

El objetivo del proyecto es desarrollar un sistema de información en forma de aplicación web que muestra la monitorización en tiempo real de alertas. \newline

Estas alertas pueden ser de tres tipos:
\begin{itemize}
    \item \textbf{Tareas}: son alertas periódicas creadas para los diferentes pacientes.
    \item \textbf{Alarmas}: son alertas que se disparan una única vez.
    \item \textbf{Presencias}: son alertas que indican la presencia de un trabajador.
\end{itemize}

El trabajo a realizar para llevar a cabo la construcción de la aplicación web consta de dos partes.

En primer lugar, implementar un frontend completo desde cero separado de la aplicación actual de Ibernex. Y en segundo lugar, el desarrollo de un backend que se realiza integrando código en la aplicación de Ibernex aprovechando algunas de las funcionalidades existentes y desarrollando las nuevas funcionalidades necesarias. \newline

La limitación a la hora de realizar el proyecto es que al tener ya una aplicación desarrollada con ciertas tecnologías, se debe realizar el análisis de tecnologías y análisis y diseño del sistema en base a esas condiciones y teniendo en cuenta las necesidades que se podrán ver posteriormente en la sección de requisitos del sistema. 


\section{Herramientas de trabajo}

% TODO: Explicar las tecnoloías que estoy utilizando para cada una de las partes, y las librerías utilizadas.

TODO: ¿Explicar frameworks, tecnologías, entornos de desarrollo, editores de código fuente, git, tortoisegit, github, google met y chat/gmail para comunicación con empresa en remoto?

\section{Esquema general de la memoria del proyecto}

% TODO: Exponer las secciones que pertenecen a la memoria del proyecto para hacer una presentación de lo que se explica en cada una de ellas y no solo con el índice y los títulos.

TODO: redactar una vez clara la estructura final de la memoria








